% Shared implementation-grounded technical specification for the ConfigStream reports.
% It is intentionally a local report source: every design statement below is tied to
% the implementation under processing/, spark/, api/, and the registered workloads.

\appendix
\section{Technical Specification of the ConfigStream Pipeline}
\label{app:technical-specification}

This appendix turns the architectural overview into an implementable specification. Its purpose is twofold. First, it records precisely what the current system does, so that a reader can distinguish an implemented property from an aspiration. Second, it explains why several apparently simple choices---for example, comparing a new Kafka row with the preceding row---would change the meaning of a configuration history. The specification describes the currently supported inline-text configuration path. It does not claim a general parser for all configuration formats, a transactional protocol across all sinks, or a high-availability deployment.

\subsection{Problem setting and terminology}

ConfigStream observes versions of a configuration-bearing asset. An asset is a durable logical subject, such as a router, host, application, or managed resource. A snapshot is one observed configuration associated with that asset. A producer version is the ordering label supplied for the asset history. An arrival is a delivery to Kafka. These terms are deliberately distinct: one snapshot may be redelivered; two deliveries may be reversed; and an arrival time is not evidence that one configuration semantically follows another.

For one asset $a$, let $S_a$ denote the set of accepted snapshots and let $v(s)$ denote the producer version of snapshot $s$. The desired historical relation is the set of justified adjacent edges
\[
E_a = \{(s_i,s_j) : v(s_j)=v(s_i)+1,\; s_i,s_j\in S_a\}.
\]
An edge is therefore a statement about two accepted versions, not an observation about two neighboring transport records. The definition intentionally leaves a gap visible. If versions 18 and 20 are present but version 19 is absent, then neither the existence nor the content of an 18-to-20 transition establishes an 18-to-20 adjacent change. The offline comparison utility can compare explicitly nominated endpoints and report the resulting version gap; the online stateful path does not silently reinterpret that comparison as an adjacent transition.

This distinction has practical consequences for incident analysis. A dashboard answer to ``what was current?'' is based on the maximum accepted producer version for the asset. An answer to ``what changed?'' is based on an adjacent edge. An answer to ``why is a change absent?'' can refer to a diagnostic such as a gap, conflict, state eviction, or state timeout. Treating all three answers as a single arrival-ordered log query would make them concise but unreliable.

\subsection{Typed records and identity contract}

The codebase keeps the public records separate because they carry different provenance and recovery obligations. The snapshot contract contains the stable asset identifier, immutable snapshot identifier, producer version, event time, ingest time, content or content URI, raw SHA-256, and schema version. A normalized snapshot adds policy-dependent normalized content and normalized hash. A diff identifies two normalized endpoints and records line-level change information. An alert is a rule evaluation over a valid diff. A diagnostic records a safe failure category rather than treating rejected input as ordinary data.

Table~\ref{tab:contract-roles} summarizes the role of the main contracts. The table is not merely a schema inventory. It tells the reader which fields may be used to make a decision and which fields are retained only as provenance. In particular, a raw hash answers whether source bytes match the supplied source representation; a normalized hash answers whether the supported normalization policy produced the same canonical text. Neither is an ordering key.

\begin{table}[ht]
\centering
\caption{Principal record contracts and the decisions they support.}
\label{tab:contract-roles}
\small
\begin{tabularx}{\textwidth}{p{0.23\textwidth}X X}
\toprule
Contract & Required role in the pipeline & Deliberate non-claim \\
\midrule
Asset & Stable logical subject through \texttt{asset\_id}. & It does not prove that a real-world device was reachable. \\
ConfigSnapshot & Versioned source observation with source hash and event-time provenance. & It does not prove global uniqueness or cross-topic ordering. \\
NormalizedConfig & Canonical text and normalized hash under one named policy. & It does not preserve every source formatting distinction. \\
ConfigDiff & Deterministic comparison of two admitted endpoints. & It does not imply that non-adjacent versions are a valid transition. \\
Alert & Rule outcome attached to a diff, rule identifier, and new event time. & It is not an exactly-once external notification. \\
StateDiagnostic & Explicit record of conflict, eviction, timeout, or unresolved gap. & It is not evidence that the missing configuration was deleted. \\
Lake manifest & Batch-level inventory, counts, paths, and file hashes. & It is not a distributed transaction with search and Kafka. \\
\bottomrule
\end{tabularx}
\end{table}

The first admission boundary is identity. The identity processor binds a \texttt{snapshot\_id} before version-state processing. An exact retry of a retained identity is suppressed. A reuse of the same snapshot identifier with incompatible identity fields is emitted as \texttt{snapshot\_identity\_conflict}. This is stronger than deduplicating arbitrary payload bytes: it makes the source's immutable identifier part of the replay contract. It is still bounded by retained state and configured timeout, so it is not a claim that an unbounded global duplicate ledger exists.

The second boundary is version compatibility within an asset. If an accepted buffer already contains version $k$ with a different snapshot identifier, a later record at $k$ yields \texttt{version\_conflict}. If a candidate has a version at or below the maintained eviction floor, it yields \texttt{evicted\_version}. These outcomes are intentionally observable. Suppressing them without a diagnostic would make the resulting history look complete even though the processor had lost the information required to judge an old arrival.

\subsection{Ingestion, validation, and event-time interpretation}

The online query consumes \texttt{config.snapshots}. Kafka provides a partitioned replay log, but the report does not equate its partition order with asset history. The source is configured with a maximum offsets-per-trigger bound, and validation happens before an input participates in state. The validation layer rejects records with unsupported schema version, missing identifiers, invalid non-positive version, malformed timestamps, or inconsistent source metadata. Invalid records are separated into safe diagnostics; configuration bodies are not printed as an error-reporting convenience.

For valid snapshots, Structured Streaming assigns the declared event-time column and applies the configured watermark delay. The default configuration sets the delay to ten minutes. The watermark is a resource and lateness policy, not a declaration that all source clocks are correct. A snapshot inside the retained event-time horizon can still complete an earlier missing edge. A snapshot that arrives after the relevant state has expired cannot be reconstructed from information no longer held by the state store. The report consequently describes the system as boundedly repairable under disorder, not as losslessly reorderable for arbitrary lateness.

The implementation has two stateful stages. The first groups source identities to implement retained duplicate and conflict handling. The second groups valid normalized candidates by \texttt{asset\_id} and operates over producer versions. Keeping these stages separate makes the failure model easier to reason about: identity reuse is a source contract issue; adjacency is a temporal-history issue. Combining them into one undocumented key would make it harder to establish whether a rejected record was a duplicate, a version collision, or a late arrival beyond retention.

\subsection{Version-buffer state machine}

The pure transition model in \texttt{processing/state.py} is the semantic reference for the Spark implementation. Its state is a tuple of buffered snapshots sorted by version, a set of predecessor versions whose next edge has already been emitted, and an eviction floor. Each buffered entry retains the normalized configuration, snapshot identifier, producer version, and event time in microseconds. The Spark \texttt{VersionProcessor} persists the corresponding records in list state \texttt{versions\_v1} and stores the floor in value state \texttt{eviction\_floor\_v1}. The deliberately named state variables make the implementation's durable state visible to a reader and to state-store inspection tools.

Algorithm~\ref{alg:state-transition} states the behavior in operational terms. It is useful because it separates three actions that are frequently conflated: accepting a source record, constructing an edge, and publishing a derived output. A record can be accepted without immediately yielding a diff. Conversely, an edge can be constructed later when an earlier missing version arrives. This is exactly the intended behavior for the arrival sequence v18, v20, v19.

\begin{quote}
\textbf{Algorithm 1: Bounded adjacent-edge transition.} Given a state $(B,F,M)$ and a candidate $x$: (1) reject $x$ if its retained identity conflicts, if its version collides with another identity, or if $v(x)\leq F$; (2) otherwise insert $x$ into $B$ and sort by version; (3) for every neighboring pair $(p,q)$ in $B$, emit an edge only when $v(q)=v(p)+1$ and $v(p)\notin M$; then add $v(p)$ to $M$; (4) if $|B|$ exceeds the cap, emit any newly justified edges, evict the lowest-version record, set $F$ to that version, and emit an eviction or unresolved-gap diagnostic as applicable; (5) retain only edge flags whose two endpoints remain retained.\label{alg:state-transition}
\end{quote}

The default cap is 32 versions per asset. This cap is not a magic performance value; it is a bounded-memory policy. On overflow, the oldest retained version is evicted after the processor has emitted every newly justified edge available at that point. The eviction floor prevents a later arrival of an evicted old version from being treated as a fresh predecessor. Without the floor, a delayed v18 could be reinserted after newer history had been evicted and could create a misleading edge against a different retained context.

The processor also registers an event-time inactivity timer. The default timeout is 30 minutes. When the timer fires after the watermark reaches the asset's expiry time, the state transition emits diagnostics for unresolved internal gaps and a \texttt{state\_timeout} outcome, then clears state. The timeout has an epistemic function as well as an operational one: it says that the processor no longer has enough retained context to continue repairing this asset's sequence. It does not say that the real asset stopped changing.

\subsection{Worked temporal examples}

The following examples define the expected behavior more clearly than a generic statement that the system ``handles out-of-order data.'' Consider three snapshots of the same asset with versions 18, 19, and 20. Suppose the first micro-batch contains v18 and v20. After admission, the buffer is [18,20]. There is no consecutive pair, so no diff is emitted. If v19 then arrives before expiry, the sorted buffer becomes [18,19,20]. The processor emits the two logical edges 18--19 and 19--20 exactly once. The current/latest accepted snapshot is v20 even though v19 was the final delivery.

Now consider a retry of v19 with the same snapshot identity. The identity stage recognizes the retained retry and suppresses it; no edge is repeated. If the retry instead reuses version 19 with another snapshot identifier, it is a version conflict and does not replace the retained v19. If an asset receives versions 18 and 20 but v19 never arrives, the expiry path emits an unresolved-gap diagnostic rather than generating 18--20. Finally, if enough newer records cause v18 to be evicted, an eventual arrival of v18 is diagnosed as evicted history rather than used to create a retrospective edge.

These cases establish the local invariant used throughout the report: every emitted streaming diff corresponds to two retained, admitted snapshots of one asset whose producer versions differ by exactly one. The converse is also bounded: every such pair is emitted provided its endpoints coexist in retained state and its predecessor edge flag has not already been set. The wording matters because unbounded completeness would be false once an asset exceeds its configured capacity or timeout.

\subsection{Normalization and deterministic comparison}

Adjacent endpoints are not compared as unexamined raw byte strings. The normalizer first rejects unsupported representations, unavailable inline content, and unsafe multiline secret syntax. For supported inline text, it converts CRLF and CR line endings to LF, removes trailing whitespace, removes selected generated or last-modified comments, canonicalizes selected separators, and replaces recognized secret values with a fixed mask. The policy is intentionally conservative: it makes a small, auditable family of representational differences invisible rather than pretending to understand every vendor grammar.

Normalization produces a second representation with a separate hash. This distinction protects interpretation. A raw-hash difference with no normalized difference can mean that source bytes changed in a way intentionally ignored by the policy, such as a generated timestamp or trailing spaces. A normalized-hash difference means the selected canonical text differs. Neither result proves the operational effect of a configuration command; semantic equivalence beyond the supported policy remains outside the claim.

The comparator validates that the endpoints refer to the same asset, that the current version is strictly greater than the previous version, and that snapshot identifiers differ. It then applies line-oriented \texttt{SequenceMatcher} logic with \texttt{autojunk=False}. The fixed setting avoids data-size-dependent heuristic suppression changing the diff solely because a configuration is large or repetitive. The comparison produces additions, removals, and replacement-aligned changed lines. The latter is an overlap descriptor, not an independent count to be added to additions and removals.

The diff identifier is deterministic over the policy name, asset identifier, endpoint snapshot identifiers, endpoint versions, and normalized hashes. Consequently, a replay of the same logical pair can be recognized by its identity even if it is physically delivered again. Determinism supports validation and serving deduplication; it does not by itself create transactional delivery to every output sink.

\subsection{Rule evaluation and alert semantics}

Rules are loaded from YAML and validated at application start. A rule evaluates normalized diff lines after a valid adjacent diff has been created. It therefore cannot convert an invalid pair or a state-gap diagnostic into a normal alert. The alert identifier is derived from the diff identity and rule identifier, while alert time follows the new snapshot's event time. This choice preserves the distinction between source-time meaning and later processing or delivery time.

The rule boundary is deliberately narrow. A line-level rule can express policy checks over the supported normalized text, such as detection of a configured risky pattern. It cannot prove reachability, vendor-specific semantic safety, or deployment success. Secret masking also means a change that only modifies a masked value may be invisible in normalized lines. The raw record remains part of protected history and provenance, but ordinary API views are designed not to expose configuration bodies by default.

\subsection{Lake-first publication protocol}

The stateful query can emit accepted snapshots, normalized snapshots, diffs, alerts, diagnostics, and a raw snapshot projection. Publication occurs through a \texttt{foreachBatch} path. The publisher persists the micro-batch, writes the lake outputs first, then performs index publication and alert delivery, and finally releases the cached batch. This order establishes the intended recovery boundary: a committed Parquet batch is the authoritative historical representation from which metadata projections can be rebuilt.

For each non-empty dataset in a batch, the lake writer writes Snappy Parquet under a path scoped by query identifier and batch identifier. It inventories the resulting Parquet files, recording each path, byte length, and SHA-256. It rereads the output under the expected schema and checks the physical row count. Only after these checks does it write a versioned manifest with the query identifier, batch identifier, dataset counts, dataset locations, and file inventory. Readers use the manifest to identify committed content rather than assuming any partially created directory is complete.

The replay path is equally important. If a batch identifier already has a manifest, the writer validates the manifest and compares the expected logical projection with the committed projection. If they agree, the existing Parquet output is retained rather than overwritten. If they disagree, the conflict is surfaced instead of silently accepting an ambiguous recovery. This is a meaningful local idempotence mechanism for the lake boundary.

The protocol is not an atomic transaction across MinIO, Elasticsearch, and Kafka. A search write can fail after lake commit; a Kafka alert retry can create a physical duplicate. The report therefore claims that search documents and API views are rebuildable from committed lake data, and that deterministic identifiers allow logical deduplication where a sink supports it. It explicitly does not claim end-to-end exactly-once alert delivery.

\subsection{Serving boundary and API behavior}

The serving layer is a read-only metadata projection. Elasticsearch is used for searchable document and aggregate views, while FastAPI exposes list, lookup, latest-snapshot, alert-statistics, and diff-statistics operations. The API's latest-snapshot lookup sorts by producer version descending and snapshot identifier ascending. That tie-breaker makes the read result deterministic and preserves the principle already used by the state machine: ``latest'' is a version-semantic query, not a request for the row most recently indexed.

The API does not make the search index the authoritative history. A transient index failure after a lake commit leaves an incomplete projection that can be repaired by rebuilding from Parquet. This separation is especially valuable when debugging an apparent API contradiction: one can inspect the manifest and committed dataset before deciding whether the fault lies in state construction, lake publication, indexing, or query logic. It also bounds the sensitive-data surface because normal list and aggregate endpoints can return identifiers and metadata without returning configuration bodies.

\subsection{Evaluation model and acceptance criteria}

The evaluation is designed around evidence roles rather than a single headline throughput number. Versioned configuration history, canonical RPKI members, DNS measurements, and RIS peer-prefix occurrences have different units and different completeness assumptions. A file/s rate, a canonical-members/s rate, and an occurrences/s rate are therefore reported separately. Combining them would make a larger number but would not measure one common object.

For correctness, the primary oracle is exact identity-to-full-record projection comparison where an independent projection is available. The evaluator checks missing identities, unexpected identities, conflicting identities, duplicate identities, and full record equality. Counts and hashes are supporting checks, not substitutes for membership comparison. State tests cover out-of-order input, exact retry, identity conflict, version conflict, missing versions, cap-triggered eviction, and expiry. Lake tests cover schema reads, physical counts, manifest inventories, and replay behavior. API tests cover the version-descending latest ordering independently from ingest order.

For performance, the report states the boundary before reporting a time. The matched normalized-core campaign compares B0, B2, and B3 on one fixed complete output projection. B0 is the direct local Arrow-based rewrite, B2 is the Spark batch rewrite, and B3 is the queued Spark streaming rewrite. Parser construction, source acquisition, and production-specific group finalization are excluded in every mode. Each successful trial must exactly match the independent normalized output projection; a fast run without correct output is rejected. The accepted campaign contains three fresh-process repetitions for every workload/mode combination, yielding 27 successful trials.

The TerraDS integration replay measures a wider end-to-end boundary: Kafka publication, stateful processing, lake writing, and drain. It processed 10,188 snapshots, constructed 3,683 adjacent diffs, committed 29,402,997 bytes of raw and processed Parquet, and completed the measured drain in 68.742 seconds with zero emitted alerts and diagnostics in the accepted artifact. The corresponding publisher rate and accepted logical rate have different denominators and are therefore retained separately. These observations support an implementation-level replay claim in the declared local environment; they do not establish distributed capacity or high availability.

\subsection{Threats to validity and explicit non-claims}

The technical design is deliberately evidence-bounded. The development topology uses a single local Kafka broker with replication factor one and loopback plaintext. Many acceptance and benchmark runs execute Spark locally. The repository does not provide a completed multi-node failure campaign, cloud-object-store evaluation, or an all-sink transaction. The appropriate conclusion is that the implementation makes its state, commit, and recovery boundaries inspectable under the tested scope, not that it has demonstrated production-grade HA.

The workload sources also carry different limitations. Git history is an observable proxy for configuration evolution, not a direct measure of deployment frequency. RPKI is a structured authorization dataset; DNS observations include measurement and absence uncertainty; RIS contains collector occurrences rather than a complete global routing state. Accelerated replay rates represent a test load over accumulated history, not a natural per-device change rate. These distinctions are retained in the report because removing them would make comparisons easier to read but less defensible.

Finally, no component should be credited with a property supplied by another component. Kafka makes a replay log available; it does not resolve configuration adjacency. Spark provides managed state and event-time mechanisms; it does not define source identity. Parquet and manifests provide a durable reconstruction boundary; they do not make search writes atomic. Rules make selected normalized patterns observable; they do not provide universal semantic validation. The contribution of ConfigStream is the explicit composition of these boundaries into a testable configuration-history pipeline.

\section{Traceability Matrix from Claims to Code and Evidence}
\label{app:traceability}

Table~\ref{tab:traceability} is included to make the report auditable. A claim is only as strong as the code path and evidence that support it. The entries intentionally distinguish implemented mechanisms from evaluation artifacts and from unmeasured claims.

\begin{longtable}{p{0.26\textwidth}p{0.31\textwidth}p{0.35\textwidth}}
\caption{Traceability from report claim to implementation and evidence boundary.}\label{tab:traceability}\\
\toprule
Claim & Implementation locus & Evidence interpretation \\
\midrule
\endfirsthead
\toprule
Claim & Implementation locus & Evidence interpretation \\
\midrule
\endhead
Adjacent streaming edges require consecutive producer versions. & \texttt{processing/state.py:accept\_snapshot}; \texttt{spark/streaming/stateful.py:VersionProcessor}. & Unit and Spark state tests; out-of-order v18/v20/v19 example. Bounded by state retention. \\
Latest snapshot is chosen by version rather than arrival order. & Stateful buffer sort; \texttt{api/service.py:latest\_snapshot}. & API tests order version descending, then snapshot identifier. \\
Identity reuse and version collisions are visible. & \texttt{inspect\_snapshot}; \texttt{IdentityProcessor}. & Diagnostics are output records; they are not silently treated as accepted state. \\
Normalization is deterministic for supported inline text. & \texttt{processing/normalization.py}. & Tested policy behavior; unsupported and unsafe syntax is rejected rather than generalized. \\
Diff identity is stable for one logical pair. & \texttt{processing/diff.py}. & Identifier is derived from policy, endpoints, versions, and normalized hashes. \\
Lake output has a committed inventory boundary. & \texttt{spark/streaming/sinks/lake.py}. & Manifest, file hashes, schema reread, and row-count checks validate one committed batch. \\
Search/API projection is rebuildable. & Lake-first publisher plus search/API services. & Rebuildability follows from authoritative committed Parquet; it is not atomic cross-sink publication. \\
Matched B0/B2/B3 results are exact-output comparisons. & Registered workload evaluators and retained Phase 20 artifact. & 27 accepted corrected trials under the declared normalized-core boundary. \\
End-to-end TerraDS replay completed locally. & Kafka, streaming state, lake sink, and replay harness. & Accepted artifact reports 10,188 snapshots and 3,683 diffs; no cluster-scale inference. \\
Exactly-once delivery or HA. & No all-sink transaction or multi-node replica protocol in scope. & Not claimed; requires additional experiments and deployment design. \\
\bottomrule
\end{longtable}

The matrix also indicates how the report should be maintained. A future change to the state transition requires revisiting the adjacency claim, the out-of-order tests, the figure, and the evaluation oracle. A new output sink requires a new publication-boundary analysis rather than inheriting the lake's recovery semantics. A new parser or normalization policy requires a new policy name and an explicit account of what it preserves. In this sense, traceability is part of the methodology rather than documentation written after implementation.
